% !TEX root = ../notes.tex

\documentclass[../notes.tex]{subfiles}

\begin{document}

\section{April 13}
Here we go.

\subsection{Animated Rings}
It turns out that we do have say something concrete about animation.
\begin{definition}[animated ring]
	An \textit{animated ring} is a (finite) product-preserving functor $\mathrm{Poly}\opp\to\mathrm{Spaces}$, where $\mathrm{Poly}\subseteq\mathrm{Ring}$ is the full subcategory of polynomial $\ZZ$-algebras. Notably, $\mathrm{Poly}$ is a classical category.
\end{definition}
Let's motivate this definition. It is easier to understand the category of colimit-preserving functors.
\begin{remark}
	For any category $\mc C$, functors $\mc C\opp\to\mathrm{Spaces}$ form the category obtained by freely adding colimits to $\mc C$. (This is some version of the Yoneda embedding.) Thus, a colimit-preserving functor out of $\op{Fun}(\mc C\opp,\mathrm{Spaces})$ merely amounts to the data of a functor out of $\mc C$.
\end{remark}
Analogously, the category of product-preserving functor $\mc C\opp\to\mathrm{Spaces}$ produces the category given by added sifted colimits to $\mc C$.
\begin{definition}[sifted colimit]
	A functor $F$ preserves \textit{sifted colimit} if and only if it preserves filtered colimits and colimits indexed by $\Delta$.
\end{definition}
\begin{example}
	The moral is that an animated ring is basically a sifted colimit of $\mathrm{Poly}$. Accordingly, an animated ring is basically a colimit over $\Delta$ of some filtered colimit of polynomial $\ZZ$-algebras. Thus, it is basically a simplicial $\ZZ$-algebra.
\end{example}
\begin{remark}
	The image of the Yoneda embedding turns out to provide a set of compact projective generators. Compactness is not hard, but projectivity requires some argument (and a definition of projectivity).
\end{remark}
Here are some examples.
\begin{example}
	Let $R$ be a discrete ring, and let $\mc C$ be the classical category of free finitely generated $R$-modules. Then the category of product-preserving functors $\mc C\opp\to\mathrm{Spaces}$ is equivalent to the category of connective $R$-modules. This is a version of the Dold--Kan correspondence, where we basically send a simplicial $R$-module to the corresponding complex.
\end{example}
\begin{example}
	Let $\mc C$ be the ordinary category of finite sets. Then the category of product-preserving functors $\mc C\opp\to\mathrm{Spaces}$ is equivalent to $\mathrm{Spaces}$. The idea is that we can determine such a functor by where it sends a singleton.
\end{example}
\begin{example}
	Let $\mc C$ be the full subcategory of $\mathbb E_\infty$-rings of the form $\op{Free}_{\mathbb E_\infty}\mathbb S^{\oplus n}$. Then product-pre\-serving functors $\mc C\opp\to\mathrm{Spaces}$ is equivalent to the category of connective $\mathbb E_\infty$-rings.
\end{example}
With any kind of ring, we are interested in its modules.
\begin{remark}
	There is a functor $U$ from animated rings to free $\mathbb E_\infty$-$\ZZ$-algebras which preserves sifted colimits and sends finitely generated polynomial $\ZZ$-algebras to the corresponding $\mathbb E_\infty$-rings. This functor $U$ can be thought of as a forgetful functor; for example, for a discrete ring $R$, its corresponding animated ring goes to $R$ again under $U$. In more fancy terms, $U$ makes the category of animated rings monadic over connective $\mathbb E_\infty$-$\ZZ$-algebras. Also, $U$ preserves tensor products.
\end{remark}
\begin{definition}[module]
	Fix an animated ring $R$. Then we define $\mathrm{Mod}(R)\coloneqq\mathrm{Mod}(U(R))$.
\end{definition}
\begin{remark}
	Because $R$ has ``more structure'' than $U(R)$, one may expect $\mathrm{Mod}(R)$ to have more structure than one would a priori expect. This is true: roughly speaking, the tensor product not only gives a symmetric monoidal structure but also a $\ZZ$-linear structure.
\end{remark}
\begin{remark}
	Here is a by-product of our discussion: if $R$ is animated, then $R/x$ is animated for any $x\in\pi_0R$. Indeed, $R/x=R\otimes_{\ZZ[x]}\ZZ$ because the cofiber $R/x$ is $R\otimes_{\ZZ[x]}(x\colon\ZZ[x]\to\ZZ[x])$. Then the tensor product continues to be animated, roughly speaking because the forgetful and free functors commute with this colimit and the tensor product.
\end{remark}
Recall that the previous remark very much does not hold in the larger class of $\mathbb E_\infty$-rings.

\subsection{The Algebraic Cotangent Complex}
We can also do some deformations.
\begin{example}
	Fix a module $M$ over an animated ring $A$. Then the trivial square-zero extension $A\oplus M$ can be found as an animated ring.
\end{example}
With deformations, we can define our cotangent complex by universal property.
\begin{definition}[algebraic cotangent complex]
	Fix an animated ring $A$. There is an \textit{algebraic cotangent complex $L_A^{\mathrm{alg}}$} which satisfies
	\[\op{Hom}_{\mathrm{AniRing}_A}(A,A\oplus M)=\op{Hom}_{\mathrm{Mod}(A)}(L_A^{\mathrm{alg}},M)\]
	for any $A$-module $M$. In general, for a map $A\to B$ of animated rings, we define the \textit{relative algebraic cotangent complex} $L_{B/A}^{\mathrm{alg}}$ to satisfy
	\[\op{Hom}_{\mathrm{AniRing}_B(A)}(B,B\oplus M)=\op{Hom}_{\mathrm{Mod}(B)}(L_{B/A}^{\mathrm{alg}},M)\]
	for any $B$-module $M$.
\end{definition}
\begin{remark}
	By definition, we see that $L_A^{\mathrm{alg}}=L_{A/\ZZ}^{\mathrm{alg}}$.
\end{remark}
\begin{remark}
	It turns out that
	\[L_A^{\mathrm{alg}}=A\otimes_{A\otimes_{\mathbb S}\ZZ}L_{A/\ZZ}.\]
	(Here, $A^+\coloneqq A\otimes_{\mathbb S}\ZZ$ relates to automorphisms of the forgetful functor; notably, the $\mathbb E_\infty$-ring structure on $A^+$ depends funnily on the animated ring structure on $A$.) Thus, we can relate this cotangent complex with the one from last class. For example, $L_{B/A}^{\mathrm{alg}}=0$ if and only if $L_{B/A}=0$. One can also use this to show that $L_{A/\ZZ}^{\mathrm{alg}}$ being almost perfect implies that $L_{A/\ZZ}$ is almost perfect.
\end{remark}
The moral is that we can check facts about the cotangent complex by passing to the algebraic cotangent complex. Thus, we are motivated to do some calculations with the algebraic cotangent complex.
\begin{example}
	Fix a discrete ring $A$, and define the discrete ring $B\coloneqq A[x_1,\ldots,x_n]$. Then it turns out that $L_{B/A}^{\mathrm{alg}}$ is the free $B$-module generated by the symbols $dx_i$.
\end{example}
\begin{example}
	Fix a discrete ring $A$. Then for any $A$-alegbra $B$, one finds that $L_{B/A}^{\mathrm{alg}}$ is given by ``derived'' K\"ahler differentials. Indeed, $L_{B/A}$ commutes with sifted colimits in animated $A$-algebras $B$, so one can recover $L_{B/A}$ via a sifted colimit.
\end{example}
\begin{example}
	Fix a perfect, discrete $\FF_p$-algebra $B$. Then $L_{B/\FF_p}^{\mathrm{alg}}=0$, so $L_{B/\FF_p}=0$. This is where the adjective ``perfect'' enters homotopy theory. For example, take $B$ to be $\FF_p\left[t^{1/p^\infty}\right]$, which is the colimit of the diagram
	\[\FF_p[t]\to\FF_p[t]\to\FF_p[t]\to\cdots,\]
	where the internal maps are given by the Frobenius $t\mapsto t^p$. Thus, $L_{B/\FF_p}^{\mathrm{alg}}$ is the colimit over a diagram with objects $\FF_p[t]\cdot dt$, where the internal maps send $dt$ to $d\left(t^p\right)=0$.
\end{example}

\subsection{Suddenly Tilting}
For our next example, we need some notation.
\begin{notation}
	Set $\mathbb S[t]$ to be the $\mathbb E_\infty$-ring $\Sigma_+^\infty\NN$, where the $\mathbb E_\infty$-structure comes because $\NN$ is a monoid. Then define $\mathbb S\left[t^{1/p^\infty}\right]$ to be the colimit of the diagram
	\[\mathbb S[t]\to\mathbb S[t]\to\mathbb S[t]\to\cdots,\]
	where the internal maps are given by the ``Frobenius'' $p\colon\NN\to\NN$.
\end{notation}
\begin{example}
	Note that $M\coloneqq L_{\mathbb S\left[t^{1/p^\infty}\right]/\mathbb S}$ is a connective module such that $M\otimes_{\mathbb S}\FF_p=L_{\FF_p\left[t^{1/p^\infty}\right]/\FF_p}=0$. Thus, we may not know much about $M$, but we do know that its $p$-completion vanishes.
\end{example}
\begin{proposition}
	Define the functor $F$ from connective $p$-complete $\mathbb E_\infty$-rings to spaces by
	\[F(R)\coloneqq\op{Hom}_{\mathbb E_\infty}\left(\mathbb S\big[t^{1/p^\infty}\big],R\right).\]
	Then $F(R)$ depends only on $\pi_0(R/p)$.
\end{proposition}
\begin{remark}
	Note that $\pi_0(R/p)=\pi_0(R)/(p)$. For example, we must show that $F$ factors through $\pi_0$.
\end{remark}
\begin{proof}[Sketch]
	Because $R$ is connective, $R=\lim_n\tau_{\le n}R$. Now, $\tau_{\le n}R\to\tau_{\le n-1}R$ is a square-zero extension. However, the cotangent complex $L_{\mathbb S\left[t^{1/p^\infty}\right]/\mathbb S}$ has vanishing $p$-completion, so the space of lifts
	% https://q.uiver.app/#q=WzAsMyxbMCwxLCJcXG1hdGhiYiBTXFxsZWZ0W3ReezEvcF5cXGluZnR5fVxccmlnaHRdIl0sWzEsMSwiXFx0YXVfe1xcbGUgbi0xfVIiXSxbMSwwLCJcXHRhdV97XFxsZSBufVIiXSxbMCwxXSxbMiwxXSxbMCwyLCIiLDIseyJzdHlsZSI6eyJib2R5Ijp7Im5hbWUiOiJkYXNoZWQifX19XV0=&macro_url=https%3A%2F%2Fraw.githubusercontent.com%2FdFoiler%2Fnotes%2Fmaster%2Fnir.tex
	\[\begin{tikzcd}[cramped]
		& {\tau_{\le n}R} \\
		{\mathbb S\left[t^{1/p^\infty}\right]} & {\tau_{\le n-1}R}
		\arrow[from=1-2, to=2-2]
		\arrow[dashed, from=2-1, to=1-2]
		\arrow[from=2-1, to=2-2]
	\end{tikzcd}\]
	is contractible. Thus, the space of maps $\mathbb S\left[t^{1/p^\infty}\right]\to R$ automatically factors through $\tau_{\le0}R=\pi_0R$. Similarly, $\pi_0R$ is a square-zero extension of $\pi_0(R/p)$, so we can descend further to $\pi_0(R/p)$.
\end{proof}
\begin{remark}
	When $R$ is discrete, one can calculate that
	\[\op{Hom}_{\mathbb E_\infty}\left(\mathbb S\big[t^{1/p^\infty}\big],R\right)\]
	is the discrete set
	\[R^\flat\coloneqq\lim\left(\pi_0(R/p)\to\pi_0(R/p)\to\cdots\right),\]
	where the internal maps are (unsurprising) the Frobenius. By construction, $R^\flat$ is a discrete perfect $\FF_p$-algebra. By extending over sifted colimits, this holds for all connective $p$-complete $\mathbb E_\infty$-rings.
\end{remark}
The previous remark grants the following definition.
\begin{definition}[Witt vectors]
	We define the functor $\mathbb W$ of \textit{Witt vectors} to take perfect discrete $\FF_p$-algebras to connective $p$-complete $\mathbb E_\infty$-rings. It is the left adjoint to the functor $\op{Hom}_{\mathbb E_\infty}\left(\mathbb S\left[t^{1/p^\infty}\right],-\right)$.
\end{definition}
\begin{example}
	By analyzing the right adjoints, one finds that $\mathbb W(A)\otimes_{\mathbb S}\FF_p$ (considered as a connective $\mathbb E_\infty$-$\FF_p$-algebra) is simply the inclusion of a discrete $\FF_p$-algebra. (Here, we note that the functor $-\otimes_{\mathbb S}\FF_p$ from connective $p$-complete $\mathbb E_\infty$-rings to connective $\mathbb E_\infty$-$\FF_p$-algebras admits the forgetful functor as an adjoint.) For example, if $A$ is a perfect $\FF_p$-algebra, we recover
	\[\mathbb W(A)\otimes_{\mathbb S}\FF_p=A.\]
\end{example}
\begin{example}
	With $A=\FF_p$, we have $\mathbb W(A)=\mathbb S^\land_p$, which has $\ZZ_p$ in its $\pi_0$. Thus, we see that $\mathbb W(A)$ has managed to deform $A$.
\end{example}

\end{document}